\section{Appendix}
\label{sec:appendix}
% Source: DECISIONS.md s1 (day-split robustness), s14 (burst splitting), s11 (guard (a) grid
% sweep), s20 + s24 (D1 definition and full grid), s25.6 (XGBoost per-seed trajectories),
% s21.6 (padding), s25.5 (full cap curve -- already in Figure~\ref{fig:f6-frontier}, main body).
% TODO: prose, and decide which tables/figures move here vs. stay in the main body once the
% page budget is known (docs/PAPER_PLAN.md "Confirm the venue, template and page limit").

\subsection{Twin fraction: sensitivity to the chosen delta}
% Source: results/phase0/cicids/twin_fraction_by_jitter.csv (delta in {0.001, 0.0015, 0.003}).
% delta=0.0015 (Figure 3's value) was chosen from the s26.2 p5 gap; this checks whether the
% picture depends on that specific choice. It does, in one respect worth stating plainly: 0.003
% does not track the FPR collapse (s26.2: already at the noise floor by jitter 0.002) as tightly
% as the two finer deltas do.
Figure~\ref{fig:f3-damage-vs-fidelity}(a) uses $\delta=0.0015$, chosen from where the
realized-distance distribution's 5th percentile jumps between jitter 0 and jitter 0.002
(Section~\ref{sec:threat-model}). Table~\ref{tab:twin-delta-sensitivity} recomputes twin
fraction at $\delta\in\{0.001, 0.003\}$, bracketing that choice by roughly 2x in each direction.
The qualitative story — a cliff within one to two jitter steps of raw fidelity, not a gradient —
holds at all three, but the exact location is not delta-invariant: $\delta=0.001$ and
$\delta=0.0015$ both collapse to near-zero by jitter 0.002, matching where the FPR damage itself
disappears (s26.2); $\delta=0.003$ is coarse enough that it still reads 0.173 at jitter 0.002 —
comparable to $\delta=0.0015$'s own raw-jitter value — and does not fall to near-zero until
jitter 0.005, one step later than the damage does. A delta chosen too loosely lags the effect it
is meant to track; the finer two values, including the one used in Figure~\ref{fig:f3-damage-vs-fidelity},
do not.

\begin{table}[t]
  \centering
  \caption{Twin fraction by jitter, at three values of $\delta$.}
  \label{tab:twin-delta-sensitivity}
  \small
  \begin{tabular}{lrrr}
    \toprule
    Jitter & $\delta=0.001$ & $\delta=0.0015$ & $\delta=0.003$ \\
    \midrule
    0.000 & 0.057 & 0.105 & 0.212 \\
    0.002 & 0.00002 & 0.003 & 0.173 \\
    0.005 & 0.000 & 0.000 & 0.001 \\
    $\ge$0.01 & 0.000 & 0.000 & 0.000 \\
    \bottomrule
  \end{tabular}
\end{table}

\subsection{Guard (a): near-duplicate grid sweep, and why \texorpdfstring{\texttt{leak\_warning}}{leak\_warning} is attack-only}
% Source: DECISIONS.md s11 (grid sweep), s2/s15 (leak_warning gating rationale, moved here from
% Section 4 for space).
The near-duplicate guard (Section~\ref{sec:methodology}) removes a large fraction of the benign
class on CICIDS2017 at any reasonable grid resolution, so that resolution cannot be picked by
feel. Sweeping \texttt{near\_dup\_grid} $\in \{\text{off}, 0.005, 0.01, 0.02, 0.05\}$: with the
guard off, partition construction fails outright (CICIDS2017 has roughly 12.3k flows that are
row-content-identical, to six decimal places, across \texttt{seed\_train} and
\texttt{trusted\_eval}). From 0.005 to 0.05, benign rows removed range over $2\times$ and
\texttt{seed\_train} shrinks from 456k to 267k rows, but RF AUROC stays in 0.98--0.996 and TPR in
0.95--0.97 throughout; \texttt{leak\_warning} is set at every grid tested. The result does not
depend on which grid in this range is used, and the guard is necessary regardless of grid: this
paper uses 0.005 throughout.

\texttt{leak\_warning} itself is gated on the attack class alone, not benign. Guard (c)'s purpose
in setting that flag is to catch \texttt{S0} passing its checkpoint by memorization of a specific
held-out row rather than by learning a boundary; two different benign flows, unrelated to any
attack, can legitimately carry identical flow statistics (the same web request shape occurring
twice), so a near-zero benign nearest-neighbour distance is not by itself evidence of a leak the
way the same distance on the attack class would be. This does not mean the benign side is
ignored: guard (c) measures it on the same scale and reports it unconditionally, and
Section~\ref{sec:degeneracy} is built entirely on that benign-side measurement, because it bears
directly on A1 (whose poison and whose measured damage are both benign-side) and on the
threshold-calibration fix (Section~\ref{sec:methodology}).

\subsection{Burst splitting}
% Source: DECISIONS.md s14. Erratum embedded in s14 itself: CICIDS2017's MachineLearningCVE CSVs
% carry no Timestamp column, so every "gap"/"burst" here is a file-order run length, not a
% wall-clock inter-arrival time (same caveat as Section 4's "the CICIDS split is not temporal").
Section~\ref{sec:methodology}'s guard (b) is a row-level cut, which could in principle slice
through the middle of a contiguous burst of highly similar flows (a flood tool's output, most
obviously). Burst-level splitting tests this directly: segment each (day, family) stream into
contiguous episodes by inter-flow gap ($>2$s, chosen from the measured gap distribution) and
assign each burst whole to one partition, never dividing one across a boundary
(\texttt{tests/test\_partition.py} asserts this invariant on the reconstructed union of all four
partitions). Two family groups emerge: session-like families (Bot, SSH-/FTP-Patator, DoS
GoldenEye, Web Attack variants) have a median inter-flow gap of 4--38s and thousands of bursts
per day, where a row-level cut genuinely could split one; continuous floods (DoS Hulk, DDoS, DoS
Slowhttptest, PortScan) have a median gap under 1s and essentially do not pause — DDoS is
$\sim$128k rows in 22 bursts for the entire day.

With burst splitting on, the guard-(a) grid sweep repeats at the same order of magnitude as the
row-level cut ($\text{nn\_attack\_p5}$ 0.0010--0.0058 across grid 0.005--0.05, versus 0.0010--0.0067
row-level, Section~\ref{sec:degeneracy}); \texttt{leak\_warning} is set at every grid, unchanged.
The per-family nearest-neighbour breakdown explains why burst-awareness does not help: DoS
Hulk's nearest \texttt{honeypot\_pool} neighbour to a \texttt{trusted\_eval} row has RMS distance
0.0003, DDoS 0.00017, DoS GoldenEye 0.00017 — a different burst, sometimes a different half of
the day, and the feature vector is still within noise of zero distance. The degeneracy
(Section~\ref{sec:degeneracy}) is a property of the feature representation these families
produce, not an artifact of where or how the cut is drawn.

\subsection{D1: full definition and grid}
% TODO. Source: DECISIONS.md s20, s24.

\subsection{XGBoost per-seed trajectories}
% Source: DECISIONS.md s25.6. Full 8-round, both-model per-seed table; Section 6 summarizes
% this and refers here rather than reproducing all 10 rows in the main body.
Section~\ref{sec:attack} reports XGBoost's fixed-threshold FPR response to A1 as seed-dependent
rather than characterisable from a mean over 5 seeds; Table~\ref{tab:xgb-seeds-full} is the full
per-seed trajectory this is based on (ratio 0.5, raw fidelity, rounds 2 through 20). XGBoost is
bimodal: seeds 2 and 4 reach 70--84\% FPR by round 8--10; seeds 1 and 5 stay below 2\% through
round 10, and only seed 5 climbs at all by round 20 (to 12\%); seed 3 sits in between, jumping
late (round 6--8). RF climbs on all 5 seeds, with none staying flat. A mean-over-seeds summary
hides this: XGBoost's round-10 mean is dragged up by two seeds while three show no effect at
all, which is a bimodal outcome, not a noisy estimate of one common trajectory. Every mean
XGBoost number elsewhere in this paper should be read with that in mind.

\begin{table}[t]
  \centering
  \caption{Fixed-threshold FPR by seed and round, ratio 0.5, raw fidelity, both models.}
  % Source: DECISIONS.md s25.6 (results/phase0/loop/cicids_recal/jobs/a1_{model}_s*_j0.0_r0.5_vt0.1.csv).
  \label{tab:xgb-seeds-full}
  \small
  \begin{tabular}{llrrrrrrrr}
    \toprule
    Model & Seed & r2 & r4 & r6 & r8 & r10 & r12 & r15 & r20 \\
    \midrule
    XGBoost & 1 & 0.005 & 0.005 & 0.004 & 0.006 & 0.006 & 0.006 & 0.007 & 0.019 \\
    XGBoost & 2 & 0.016 & 0.503 & 0.737 & 0.782 & 0.823 & 0.878 & 0.924 & 0.955 \\
    XGBoost & 3 & 0.006 & 0.007 & 0.053 & 0.538 & 0.693 & 0.714 & 0.734 & 0.800 \\
    XGBoost & 4 & 0.046 & 0.707 & 0.776 & 0.839 & 0.908 & 0.949 & 0.948 & 0.958 \\
    XGBoost & 5 & 0.005 & 0.004 & 0.004 & 0.006 & 0.007 & 0.008 & 0.010 & 0.117 \\
    RF & 1 & 0.048 & 0.232 & 0.618 & 0.721 & 0.767 & 0.804 & 0.858 & 0.911 \\
    RF & 2 & 0.032 & 0.153 & 0.503 & 0.700 & 0.735 & 0.776 & 0.840 & 0.894 \\
    RF & 3 & 0.021 & 0.088 & 0.362 & 0.612 & 0.706 & 0.743 & 0.776 & 0.853 \\
    RF & 4 & 0.015 & 0.095 & 0.383 & 0.615 & 0.728 & 0.740 & 0.788 & 0.869 \\
    RF & 5 & 0.042 & 0.284 & 0.647 & 0.729 & 0.779 & 0.836 & 0.880 & 0.917 \\
    \bottomrule
  \end{tabular}
\end{table}

\subsection{Cost padding}
% TODO. Source: DECISIONS.md s21.6.

\subsection{Additional limitations}
% Source: DECISIONS.md s14 erratum (file-order split), s25.5 (cap is model/dataset-specific),
% s25.7 (session-level D1/adaptive-attacker parked, not run). Moved from Section 8 for space;
% these are real limitations, not lower-priority ones -- kept complete here rather than cut.
Beyond the four limitations in Section~\ref{sec:discussion}:

\begin{itemize}
  \item \textbf{The partition is by file order, not a verified capture timestamp.}
    CICIDS2017's public \texttt{MachineLearningCVE} release carries no timestamp column
    (Section~\ref{sec:methodology}), so "day split" in this paper means the dataset's own file
    boundaries, not confirmed calendar time. A burst-aware split gives comparable results
    (Appendix~\ref{sec:appendix}, "Burst splitting"), but that is evidence the result survives a
    different file-order cut, not independent evidence of separation on real wall-clock time.
  \item \textbf{ShareCap's safe cap range is not a universal constant.} It is measured for one
    model and one dataset at a time (Section~\ref{sec:defenses}); a different model or dataset
    may have a different safe range, and we do not have a way to predict it without measuring it.
  \item \textbf{No adaptive attacker against ShareCap beyond filling the cap.} We do not test an
    attacker who observes or infers the cap and tunes its poison ratio specifically around it,
    e.g. staying just under the cap indefinitely rather than pushing past it. Session-level
    cost-of-influence weighting (a defense that might address a session-aware version of this)
    was considered and parked, not run: the dataset that would carry the source IPs and
    timestamps it needs is form-gated and was unavailable to us.
\end{itemize}
